% Input fragment; use the Chinese Volume I preamble.
% Assign a unique label before integration; remove unused optional blocks.
\chapter{待定章节标题}\label{template:v1}
\begin{chaptersummary}
【作者待填写：约 100--250 字，说明前章结果留下的问题、本章的模型或
离散选择、后续用途，以及一个核心思想。内容以正文为依据。】
\end{chaptersummary}
\learninggoals{【填写 2--4 个可检验目标。】}
% Volume I studyroute takes THREE arguments.
\studyroute{【明确略去的证明、推导或先备知识，并添加已核实的文献。】}
{【填写精确、通用的英文检索词】}
{【指定当前符号、假设、起点与可检查的展开任务。】}

\section{问题、描述层级与假设}
【说明研究目标、已有结果、保留的信息以及当前模型的假设。】

% Optional: retain only for a genuinely missing conceptual motivation.
\conceptbridge{【为什么从此前对象转向此处的模型、变量或操作？】}
{【作者待补充：80--300 字，解释需要这一步的原因与选择依据。】}
{【此前描述】$\longrightarrow$【当前描述】}

\section{数学骨架与最小例子}
【在此加入已确认的公式与推导，明确归一化、边界条件及符号约定。】
【给出一个说明机制的最小例子或极限检查。】

\section{后续接口与检查}
【说明本节输出如何被下一章或其他卷使用，并给出验证指标。】
\aiexercise{【从本章具体关系开始，指定逐步推导或可复现计算。】}
{【改变一个参数或假设，检验结论的适用范围。】}
